% GENERATED FILE. Edit canonical lessons, then run npm run generate:books.
% canonical-source-hash: fnv1a64:c3085ae3625d1856
% canonical-lessons: AR-C311-wilada, AR-C311-mawlud, AR-C311-qabila-midwife, AR-C311-anjaba, AR-C311-tasmiya

\chapter{A birth}
\label{ch:ar-a-birth}
\hlchaptermodality{311}

\begin{chapteropening}
\textbf{By the end of this chapter:} \emph{I can say childbirth, a newborn, to have a child, a midwife and a naming.}

Complete the last lesson of chapter 311: I can say childbirth, a newborn, to have a child, a midwife and a naming.
\end{chapteropening}

\section[\texorpdfstring{\ar{ولادة} (wilāda)}{ولادة (wilāda)}]{\ar{ولادة} (wilāda) — childbirth; a birth}

\label{lesson:AR-C311-wilada}



\begin{quote}
Before the new one: say the Arabic for to call out to, then the Arabic for a rescuer.
\end{quote}

\subsection*{You'll want to know: \ar{ولادة}}
\textbf{\ar{ولادة}} — \emph{wilāda} — \textquotedblleft{}childbirth; a birth\textquotedblright{}.

\begin{quote}
\ar{كانت} \ar{الولادة} \ar{في} \ar{الليل}.
\end{quote}

\begin{quote}
\emph{kānat al-wilāda fī al-layl.}
\end{quote}

\begin{quote}
The birth was at night.
\end{quote}

Birth certificate: \ar{شهادة} \ar{ولادة}; many countries say \ar{شهادة} \ar{ميلاد}.

\begin{cousinweb}
\textbf{\ar{ولادة}} (\emph{wilāda}) is built on the native Arabic root \emph{w-l-d} (Lane), \textquotedblleft{}to give birth\textquotedblright{}, the root of the taught \emph{walad}, a child, and \emph{mīlād}, a birthday.
\end{cousinweb}

\begin{grammarlens}[title={What you've built}]
One more word for telling what happened.
\end{grammarlens}

\subsection*{Your turn}
\begin{itemize}
\raggedright
  \item \emph{Say it:} \emph{wilāda}
  \item \emph{Say it:} \emph{wilāda}, once more
  \item \emph{Say it:} \emph{kānat al-wilāda fī al-layl.} — The birth was at night.
  \item \emph{Recall:} say the Arabic for to call out to, then the Arabic for a rescuer, then say \emph{wilāda} again
\end{itemize}

\begin{tcolorbox}[breakable,colback=teal!4,colframe=teal!35!black,title={Before you move on}]
What does \ar{ولادة} mean? (\textquotedblleft{}Childbirth; a birth\textquotedblright{}.) Say it once more.
\end{tcolorbox}
\section[\texorpdfstring{\ar{مولود} (mawlūd)}{مولود (mawlūd)}]{\ar{مولود} (mawlūd) — a newborn}

\label{lesson:AR-C311-mawlud}



\begin{quote}
Before the new one: say the Arabic for a rescuer, then the Arabic for childbirth.
\end{quote}

\subsection*{You'll want to know: \ar{مولود}}
\textbf{\ar{مولود}} — \emph{mawlūd} — \textquotedblleft{}a newborn\textquotedblright{}.

\begin{quote}
\ar{نام} \ar{المولود} \ar{بعد} \ar{أن} \ar{شرب} \ar{الحليب}.
\end{quote}

\begin{quote}
\emph{nāma al-mawlūd baʿda an shariba al-ḥalīb.}
\end{quote}

\begin{quote}
The newborn fell asleep after drinking his milk.
\end{quote}

Plural: \ar{مواليد}, \emph{mawālīd}.

The plural \ar{مواليد} also names everyone born in a given year.

\begin{cousinweb}
\textbf{\ar{مولود}} (\emph{mawlūd}) is built on the native Arabic root \emph{w-l-d} (Lane), \textquotedblleft{}to give birth\textquotedblright{}. \emph{Mawlūd} is its passive participle, \textquotedblleft{}one who is born\textquotedblright{}.
\end{cousinweb}

\begin{grammarlens}[title={What you've built}]
One more word for telling what happened.
\end{grammarlens}

\subsection*{Your turn}
\begin{itemize}
\raggedright
  \item \emph{Say it:} \emph{mawlūd}
  \item \emph{Say it:} \emph{mawlūd}, once more
  \item \emph{Say it:} \emph{nāma al-mawlūd baʿda an shariba al-ḥalīb.} — The newborn fell asleep after drinking his milk.
  \item \emph{Recall:} say the Arabic for a rescuer, then the Arabic for childbirth, then say \emph{mawlūd} again
\end{itemize}

\begin{tcolorbox}[breakable,colback=teal!4,colframe=teal!35!black,title={Before you move on}]
What does \ar{مولود} mean? (\textquotedblleft{}A newborn\textquotedblright{}.) Say it once more.
\end{tcolorbox}
\section[\texorpdfstring{\ar{قابلة} (qābila)}{قابلة (qābila)}]{\ar{قابلة} (qābila) — a midwife}

\label{lesson:AR-C311-qabila-midwife}



\begin{quote}
Before the new one: say the Arabic for childbirth, then the Arabic for a newborn.
\end{quote}

\subsection*{You'll want to know: \ar{قابلة}}
\textbf{\ar{قابلة}} — \emph{qābila} — \textquotedblleft{}a midwife\textquotedblright{}.

\begin{quote}
\ar{ساعدت} \ar{القابلة} \ar{الأم} \ar{في} \ar{الولادة}.
\end{quote}

\begin{quote}
\emph{sāʿadat al-qābila al-umm fī al-wilāda.}
\end{quote}

\begin{quote}
The midwife helped the mother with the birth.
\end{quote}

Not the taught verb \ar{قابل}, to meet: \ar{قابلة} is a noun here.

\begin{cousinweb}
\textbf{\ar{قابلة}} (\emph{qābila}) is built on the native Arabic root \emph{q-b-l} (Lane), \textquotedblleft{}to face; to receive\textquotedblright{}. \emph{Qābila} is the active participle, \textquotedblleft{}she who receives\textquotedblright{}: the midwife is usually explained as the woman who receives the child at its birth.
\end{cousinweb}

\begin{grammarlens}[title={What you've built}]
One more word for telling what happened.
\end{grammarlens}

\subsection*{Your turn}
\begin{itemize}
\raggedright
  \item \emph{Say it:} \emph{qābila}
  \item \emph{Say it:} \emph{qābila}, once more
  \item \emph{Say it:} \emph{sāʿadat al-qābila al-umm fī al-wilāda.} — The midwife helped the mother with the birth.
  \item \emph{Recall:} say the Arabic for childbirth, then the Arabic for a newborn, then say \emph{qābila} again
\end{itemize}

\begin{tcolorbox}[breakable,colback=teal!4,colframe=teal!35!black,title={Before you move on}]
What does \ar{قابلة} mean? (\textquotedblleft{}A midwife\textquotedblright{}.) Say it once more.
\end{tcolorbox}
\section[\texorpdfstring{\ar{أنجب} (anjaba)}{أنجب (anjaba)}]{\ar{أنجب} (anjaba) — to have a child, to bear}

\label{lesson:AR-C311-anjaba}



\begin{quote}
Before the new one: say the Arabic for a newborn, then the Arabic for a midwife.
\end{quote}

\subsection*{You'll want to know: \ar{أنجب}}
\textbf{\ar{أنجب}} — \emph{anjaba} — \textquotedblleft{}to have a child, to bear\textquotedblright{}.

\begin{quote}
\ar{أنجبت} \ar{جارتنا} \ar{بنتا} \ar{جميلة}.
\end{quote}

\begin{quote}
\emph{anjabat jāratunā bintan jamīla.}
\end{quote}

\begin{quote}
Our neighbour had a beautiful baby girl.
\end{quote}

Present: \ar{ينجب}, \emph{yunjib}.

Of either parent: \ar{أنجب} \ar{الرجل} \ar{ثلاثة} \ar{أولاد}.

\begin{cousinweb}
\textbf{\ar{أنجب}} (\emph{anjaba}) is built on the native Arabic root \emph{n-j-b} (Lane), the root of \emph{najīb}, noble, excellent. Form IV \emph{anjaba} was first \textquotedblleft{}to beget a noble child\textquotedblright{}, then to have a child.
\end{cousinweb}

\begin{grammarlens}[title={What you've built}]
One more word for telling what happened.
\end{grammarlens}

\subsection*{Your turn}
\begin{itemize}
\raggedright
  \item \emph{Say it:} \emph{anjaba}
  \item \emph{Say it:} \emph{anjaba}, once more
  \item \emph{Say it:} \emph{anjabat jāratunā bintan jamīla.} — Our neighbour had a beautiful baby girl.
  \item \emph{Recall:} say the Arabic for a newborn, then the Arabic for a midwife, then say \emph{anjaba} again
\end{itemize}

\begin{tcolorbox}[breakable,colback=teal!4,colframe=teal!35!black,title={Before you move on}]
What does \ar{أنجب} mean? (\textquotedblleft{}To have a child, to bear\textquotedblright{}.) Say it once more.
\end{tcolorbox}
\section[\texorpdfstring{\ar{تسمية} (tasmiya)}{تسمية (tasmiya)}]{\ar{تسمية} (tasmiya) — a naming; a name given}

\label{lesson:AR-C311-tasmiya}



\begin{quote}
Before the new one: say the Arabic for a midwife, then the Arabic for to have a child.
\end{quote}

\subsection*{You'll want to know: \ar{تسمية}}
\textbf{\ar{تسمية}} — \emph{tasmiya} — \textquotedblleft{}a naming; a name given\textquotedblright{}.

\begin{quote}
\ar{اختار} \ar{الجد} \ar{تسمية} \ar{المولود}.
\end{quote}

\begin{quote}
\emph{ikhtāra al-jadd tasmiyat al-mawlūd.}
\end{quote}

\begin{quote}
The grandfather chose the baby’s name.
\end{quote}

The verb is \ar{سمى}, to name.

\begin{cousinweb}
\textbf{\ar{تسمية}} (\emph{tasmiya}) is built on the native Arabic root \emph{s-m-w} (Lane), \textquotedblleft{}to be high\textquotedblright{}, the root of the taught \emph{ism}, a name. \emph{Tasmiya} is the verbal noun of Form II \emph{sammā}, \textquotedblleft{}to give a name\textquotedblright{}.
\end{cousinweb}

\begin{grammarlens}[title={What you've built}]
One more word for telling what happened.
\end{grammarlens}

\subsection*{Your turn}
\begin{itemize}
\raggedright
  \item \emph{Say it:} \emph{tasmiya}
  \item \emph{Say it:} \emph{tasmiya}, once more
  \item \emph{Say it:} \emph{ikhtāra al-jadd tasmiyat al-mawlūd.} — The grandfather chose the baby’s name.
  \item \emph{Recall:} say the Arabic for a midwife, then the Arabic for to have a child, then say \emph{tasmiya} again
\end{itemize}

\begin{tcolorbox}[breakable,colback=teal!4,colframe=teal!35!black,title={Before you move on}]
What does \ar{تسمية} mean? (\textquotedblleft{}A naming; a name given\textquotedblright{}.) Say it once more.
\end{tcolorbox}
